\section{Method}
\label{sec:method}

\paragraph{Residual updates as a depth-indexed sequence.} Consider a pre-norm
Transformer with $L$ blocks acting on a hidden state $\hh_\ell \in \mathbb{R}^{T
\times d}$ for $T$ tokens of width $d$. Each block contributes an additive
residual update
\begin{equation}
\hh_{\ell+1} = \hh_\ell + \dd_\ell,
\qquad
\dd_\ell = F_\ell(\hh_\ell),
\label{eq:residual}
\end{equation}
where $F_\ell$ is the attention-plus-MLP sublayer stack. Skipping block $\ell$
means producing an estimate $\widehat{\dd}_\ell$ \emph{without} evaluating
$F_\ell$, then propagating $\widehat{\hh}_{\ell+1} = \hh_\ell +
\widehat{\dd}_\ell$. Seen this way, every layer-skipping method is a
\emph{predictor} of a missing update, and methods differ only in what they
predict.

\paragraph{Plain skipping is a zero-order predictor.} The conventional skip
operator sets
\begin{equation}
\widehat{\dd}^{\,\mathrm{skip}}_\ell = 0 ,
\label{eq:plainskip}
\end{equation}
the unique estimator that uses no information at all. That is what makes it cheap,
and it is also what makes it wasteful: by the time the model reaches layer $\ell$
it has already computed $\dd_{\ell-1}, \dd_{\ell-2}, \dots$, all of which sit in
the residual stream and all of which \cref{eq:plainskip} discards. A second, equally
fitting-free estimator is available: reuse the previous update,
$\widehat{\dd}_\ell = \dd_{\ell-1}$. and we carry this \textbf{Copy Update} baseline throughout,
because any gain \mname shows over it is a gain from \emph{fitting}, not merely
from being nonzero.

\paragraph{\mname: per-channel autoregression across depth.} We model the update
sequence as autoregressive in depth, with one coefficient per channel. Writing
$\aal_\ell \in \mathbb{R}^{d}$ for the fitted coefficient vector,
\begin{equation}
\widehat{\dd}_\ell = \aal_\ell \odot \dd_{\ell-1},
\qquad
\widehat{\hh}_{\ell+1} = \hh_\ell + \aal_\ell \odot \dd_{\ell-1},
\label{eq:diag}
\end{equation}
where $\odot$ is the elementwise product, broadcast over the $T$ token positions.
Each channel $j$ has an independent closed-form least-squares solution,
\begin{equation}
a^{\star}_{\ell,j} =
\frac{\sum_t \Delta_{\ell-1,tj}\,\Delta_{\ell,tj}}
     {\sum_t \Delta_{\ell-1,tj}^{2} + \epsilon}.
\label{eq:fit}
\end{equation}
Constraining every channel of $\aal_\ell$ to a common value gives the scalar
predictor
\begin{equation}
\widehat{\dd}_\ell = \alpha_\ell \dd_{\ell-1},
\qquad
\alpha^{\star}_\ell =
\frac{\langle \dd_{\ell-1}, \dd_\ell \rangle_F}
     {\|\dd_{\ell-1}\|_F^2 + \epsilon},
\label{eq:ar1}
\end{equation}
and setting $\aal_\ell = \bm{0}$ recovers plain skipping. \mname therefore
generalizes both. We report the scalar predictor throughout as an ablation rather
than a simplification: \Cref{sec:experiments} shows that the restriction to a
single coefficient is not a minor loss of capacity but the difference between a
method that repairs a skipped block and one that, at some layers, damages it
further.

\paragraph{Closed-form calibration, no gradients.} Both estimators are fitted by
streaming a small unlabeled calibration set through the dense model once and
accumulating the sums in \cref{eq:fit} in FP64. No hidden states are stored, no
gradients are taken, no labels are used, and recalibrating an entire model costs a
single forward pass over 16--32 sequences. Padding tokens are excluded from
the accumulators.

\paragraph{Which layers to skip.} The natural companion to \cref{eq:fit} is the held-out
fraction of update energy a predictor explains,
\begin{equation}
P_\ell = 1 - \frac{\sum \|\dd_\ell - \widehat{\dd}_\ell\|_2^2}{\sum \|\dd_\ell\|_2^2},
\label{eq:pscore}
\end{equation}
computable from the dense model before any skipping decision. We evaluate it against two rules
that instead measure the effect of skipping directly on a development set
(\Cref{sec:app-selection}). \Cref{sec:experiments} shows the choice between them matters more
than the choice of predictor, and that the cheap one is the wrong one.

\paragraph{Cost.} \mname stores a $d$-dimensional coefficient vector per skipped
block ($d = 896$ on Qwen2.5-0.5B), three orders of magnitude fewer than the $d^2$
parameters a dense boundary map between hidden states would require. Inference cost is unchanged in order: one elementwise multiply and one add,
$O(Td)$ work against the $O(Td^2)$ the block itself would have cost, on tensors
already resident in the residual stream. The only additional memory traffic is
retaining $\dd_{\ell-1}$ for one further layer. The block-compute reduction that
motivates skipping therefore survives intact.
